%===============================================================================
\subsection{Evaluation Setup, Metrics and Datasets}
\label{subsec:eval_setup_metrics}
%===============================================================================

To ensure a consistent and fair benchmark across all evaluated  models, we use \gls{mmd} to assess the similarity between predicted and true gene expression distributions.
The \gls{mmd} is calculated per type of perturbation (single or combinations thereof) and then averaged to provide an overall measure of distribution similarity.
Alongside objective metric comparison, we use Precision@10 to evaluate non-additive interactions of combinatorial perturbations (see \cref{app:metrics}).

For evaluation we use two datasets: First, the Perturb-seq single-cell RNA-seq dataset from~\citet{norman2019exploring} and as processed in the CPA study~\citep{lotfollahi2023predicting} and refer to it as the Norman-CPA dataset.
It profiles 105 gene perturbations in K562 cells, including both single and double perturbations.
In total, 284 conditions across $\sim$108,000 cells were measured, of which 131 represent unique gene–gene combinations and the remaining 153 represent single perturbations. For each perturbation, associated control cells were also profiled. 
We utilize the entire set of single-perturbations during training and evaluate the subsequent methods on the double-perturbations.

Additionally, we used the CRISPRi single-cell RNA-seq dataset curated by~\citet{replogle_combinatorial_2020} and as processed in the PerturBase study~\citep{wei_perturbase_2024}. We refer to it as the Replogle2020 dataset.
This dataset profiles 82 gene perturbations in K562 cells, including both single (44), double perturbations (37), and control.
The 82 conditions were measured across $\sim$22,740 cells.